\slajdid{09_3-s036}
\begin{frame}[label=09_3-s036]
Da bi ipak malo drugačije posmatrali, a ne da svaki put izvodimo iste jednačine za različite tipove logičkih kola, zbog njihove različite ulazne i izlazne kapacitivnosti, vratićemo se na polazne jednačine kod izvođenja u lancu invertora, odnosno na izraz za kašnjenje prvog jediničnog invertora
\[t_{p1}=0.69R_{eq1}(C_{int1}+C_{ext1})\]
Uvešćemo pojam \alert{karakteristične vremenske konstante jediničnog invertora}. U tom smislu ulazna kapacitivnost, kapacitivnosti gejtova, jediničnog invertora je
\[C_{in}=C_{OX}(W_nL_n+W_pL_p)\]
Uz pretpostavku $L_n=L_p=L$ i $W_p=2W_n=2W$
\[C_{in}=3C_{OX}WL=3W_nC_g\]
gde je $C_g=C_{OX}L_n$ normirana kapacitivnost po širini kanala tranzistora. U tom slučaju karakteristična vremenska konstanta jediničnog invertora je
\[\tau_{inv}=0.69R_{eq}C_{in}=3WC_g\times0.69R_{eq}\]
\end{frame}
